\documentclass[a4paper, serif, 10pt]{moderncv}

%% moderncv
% style and color
\moderncvstyle{banking}
\moderncvcolor{black}

%% page layout/margins
\usepackage[top=2.5cm, bottom=2.5cm, left=1.5cm, right=1.5cm]{geometry}

\nopagenumbers{}

%% Babel config for French language
% ref:
% - https://latex3.github.io/babel/guides/locale-french.html
\usepackage[french]{babel}

%% fontspec
\usepackage{fontspec}

\IfFontExistsTF{Linux Libertine}{
  \setmainfont[Ligatures={TeX, Common}, Numbers=OldStyle]{Linux Libertine}
}{}
\IfFontExistsTF{Linux Libertine O}{
  \setmainfont[Ligatures={TeX, Common}, Numbers=OldStyle]{Linux Libertine O}
}{}

%% CTeX
% CJK support, used to show CJK characters in the resume
%
% - fontset=none: disable builtin fontset but instead set the CJK font manually
% - heading=false: disable ctex heading
% - punct=kaiming: use kaiming punctuations styles for CJK
% - scheme=plain: use plain scheme, do not override \normalsize font size
% - space=auto: space settings for CJK characters
%
% ref:
% - http://ctan.mirrorcatalogs.com/language/chinese/ctex/ctex.pdf
\usepackage[UTF8, heading=false, punct=kaiming, scheme=plain, space=auto]{ctex}

\IfFontExistsTF{Noto Serif CJK SC}{
  \setCJKmainfont{Noto Serif CJK SC}
}{}
\IfFontExistsTF{Noto Sans CJK SC}{
  \setCJKsansfont{Noto Sans CJK SC}
}{}

%% line spacing
\usepackage{setspace}
\setstretch{1.125}

%% URL styling - use normal text instead of monospace
\urlstyle{same}

% Auto-underline all links
% Defer until \begin{document} so hyperref is already loaded
\AtBeginDocument{%
  \let\oldhref\href
  \renewcommand{\href}[2]{\oldhref{#1}{\underline{#2}}}%
}

%% Patch \httplink and \httpslink to handle full URLs
% The original moderncv macros always prepend http:// or https:// to URLs.
% This causes issues when the URL already has a protocol (e.g., https://example.com)
% resulting in malformed URLs like https://https://example.com.
% This patch checks if the URL already contains "://" and uses it directly if so.
% Note: We use \str_set:Nx (x-expansion) to expand macros like \@homepage first.
\ExplSyntaxOn
\str_new:N \g_yamlresume_url_str

\RenewDocumentCommand{\httpslink}{O{}m}{
  \str_set:Nx \g_yamlresume_url_str {#2}
  \str_if_in:NnTF \g_yamlresume_url_str {://}
    { \url{#2} }
    { \url{https://#2} }
}

\RenewDocumentCommand{\httplink}{O{}m}{
  \str_set:Nx \g_yamlresume_url_str {#2}
  \str_if_in:NnTF \g_yamlresume_url_str {://}
    { \url{#2} }
    { \url{http://#2} }
}
\ExplSyntaxOff

\name{Camille Moreau}{}
\title{Data Scientist spécialisée en machine learning et analyse prédictive}
\phone[mobile]{+33 6 12 34 56 78}
\email{camille.moreau@email.fr}
\homepage{https://camillemoreau.fr}

\address{12 rue de la Paix, Paris, Île-de-France, France, 75002}{}{}

\extrainfo{{\small \faLinkedin}\ \href{https://www.linkedin.com/in/camille-moreau-data}{@camille-moreau-data} {} {} {} • {} {} {}
{\small \faGithub}\ \href{https://github.com/camillemoreau}{@camillemoreau} {} {} {} • {} {} {}
{\small \faMedium}\ \href{https://medium.com/@camillemoreau}{@camillemoreau}}

\begin{document}

\maketitle

\section{Informations de base}

\cvline{}{Data Scientist avec plus de 6 ans d'expérience dans la conception, le déploiement et le suivi de modèles de machine learning en environnement production.
%
\begin{itemize}
\item Expertise en Python, SQL et visualisation de données (Power BI, Tableau)
\item Expérience en modélisation prédictive, traitement du langage naturel et détection de fraude
\item Solide connaissance des environnements cloud (AWS, GCP) et des pratiques MLOps
\item Esprit analytique, rigueur scientifique et capacité à vulgariser les résultats auprès des parties prenantes non techniques
\end{itemize}}
\section{Formation}

\cventry{sept. 2017–juin 2019}
        {Master, Data Science et Intelligence Artificielle, Score : Mention Bien}
        {Sorbonne Université}
        {\url{https://www.sorbonne-universite.fr/}}
        {}
        {\begin{itemize}
\item Diplômée avec mention Bien, spécialisation en apprentissage automatique
\item Mémoire de fin d'études sur l'utilisation du deep learning pour la détection de fraude financière
\item Projet de fin d'études réalisé en collaboration avec un grand acteur bancaire
\end{itemize}
\textbf{Cours} : Apprentissage statistique, Deep Learning, Big Data Analytics, Visualisation de données, Optimisation et recherche opérationnelle, Systèmes de recommandation, Traitement automatique du langage naturel, Éthique des données}

\cventry{sept. 2014–juin 2017}
        {Licence, Mathématiques et Informatique}
        {Université Claude Bernard Lyon 1}
        {\url{https://www.univ-lyon1.fr/}}
        {}
        {\begin{itemize}
\item Formation fondamentale en mathématiques, algorithmique et programmation
\item Projets pratiques de statistiques et d'analyse de données
\item Stage de fin de licence sur l'analyse exploratoire de données de consommation énergétique
\end{itemize}}
\section{Expérience professionnelle}

\cventry{mars 2021–Present}
        {Data Scientist Senior}
        {Nova Analytics}
        {\url{https://www.nova-analytics.fr/}}
        {}
        {\begin{itemize}
\item Conception et déploiement de modèles de machine learning pour l'optimisation de la chaîne logistique
\item Développement de pipelines de données automatisés en Python et Spark, réduisant de 30 \% le temps de préparation des données
\item Mise en place de tableaux de bord interactifs avec Power BI pour suivre les indicateurs de performance
\item Encadrement d'une équipe de trois data scientists juniors, avec revue de code et suivi des bonnes pratiques
\item Collaboration avec les équipes produit et métier pour traduire les besoins en solutions analytiques
\end{itemize}
\textbf{Mots-clés} : Python, Spark, Power BI, Machine Learning, MLOps, AWS, Management}

\cventry{juil. 2019–févr. 2021}
        {Data Scientist}
        {PayLogic}
        {\url{https://www.paylogic.fr/}}
        {}
        {\begin{itemize}
\item Développement de modèles de scoring pour la détection de fraude aux transactions bancaires
\item Amélioration de la précision du modèle de détection de 15 \% grâce à l'ingénierie de caractéristiques
\item Mise en œuvre d'algorithmes de traitement du langage naturel pour l'analyse des retours clients
\item Rédaction de rapports d'analyse et présentation des résultats aux comités de direction
\end{itemize}
\textbf{Mots-clés} : Python, Scikit-learn, SQL, Détection de fraude, NLP}
\section{Langues}

\cvline{Français}{Niveau natif ou bilingue \hfill \textbf{Mots-clés} : Langue maternelle}
\cvline{Anglais}{Niveau professionnel complet \hfill \textbf{Mots-clés} : TOEIC 960, Lecture d'articles scientifiques, Présentations internationales}
\section{Compétences}

\cvline{Machine Learning}{Expert \hfill \textbf{Mots-clés} : Python, Scikit-learn, TensorFlow, PyTorch, XGBoost, LightGBM}
\cvline{Analyse et Visualisation de Données}{Expert \hfill \textbf{Mots-clés} : Pandas, NumPy, Matplotlib, Seaborn, Power BI, Tableau, SQL}
\cvline{Big Data et Cloud}{Avancé \hfill \textbf{Mots-clés} : Apache Spark, Hadoop, AWS, GCP, Docker, Airflow}
\cvline{Statistiques}{Avancé \hfill \textbf{Mots-clés} : R, Régression, Séries temporelles, Statistiques inférentielles, Tests d'hypothèses}
\section{Récompenses}

\cventry{nov. 2023}
        {Data Innovation Summit France}
        {Prix de la meilleure innovation data}
        {}
        {}
        {Récompense attribuée pour le développement d'un modèle prédictif de maintenance industrielle ayant réduit les coûts d'arrêt de production de 20 \%.}

\cventry{oct. 2019}
        {Sorbonne Université}
        {Prix du meilleur mémoire de Master}
        {}
        {}
        {Distinction décernée pour le mémoire portant sur la détection de fraude financière par réseaux de neurones profonds.}
\section{Certificats}

\cventry{sept. 2021}
        {Amazon Web Services}
        {AWS Certified Machine Learning – Specialty}
        {\url{https://aws.amazon.com/certification/}}
        {}
        {}

\cventry{avr. 2020}
        {Coursera}
        {Deep Learning Specialization}
        {\url{https://www.coursera.org/specializations/deep-learning}}
        {}
        {}
\section{Publications}

\cventry{févr. 2022}
        {Détection de fraude financière par réseaux de neurones profonds}
        {Revue Française de Data Science}
        {\url{https://www.rfds.fr/articles/detection-fraude}}
        {}
        {Publication co-écrite présentant une approche de deep learning pour détecter des transactions frauduleuses en temps réel. L'article détaille les architectures comparées et les résultats sur un jeu de données bancaires anonymisé.}
\section{Projets}

\cventry{janv. 2023–juin 2023}
        {Moteur de recommandation de contenus culturels basé sur le filtrage collaboratif et le deep learning.}
        {Système de recommandation de contenus}
        {\url{https://github.com/camillemoreau/reco-system}}
        {}
        {\begin{itemize}
\item Conception d'un système de recommandation hybride combinant filtrage collaboratif et embeddings
\item Amélioration de 12 \% du taux de clics sur une plateforme de streaming culturel
\item APIs Dockerisées déployées sur AWS Elastic Beanstalk
\end{itemize}
\textbf{Mots-clés} : Recommandation, TensorFlow, AWS, API}

\cventry{sept. 2022–déc. 2022}
        {Outil d'analyse d'opinions en temps réel à partir de Twitter.}
        {Analyse de sentiments des tweets}
        {\url{https://github.com/camillemoreau/sentiment-twitter}}
        {}
        {\begin{itemize}
\item Modèle de NLP basé sur BERT pour classer les tweets en sentiments positifs, négatifs ou neutres
\item Pipeline de collecte de données avec Twitter API et Kafka
\item Visualisation des tendances avec Grafana
\end{itemize}
\textbf{Mots-clés} : NLP, BERT, Kafka, Grafana}
\section{Centres d'intérêt}

\cvline{Randonnée}{Alpes, Mont Blanc}
\cvline{Photographie}{Photo de rue, Voyages}
\cvline{Lecture}{Science-fiction, Essais}
\section{Bénévolat}

\cventry{janv. 2022–Present}
        {Data Scientist bénévole}
        {Data for Good France}
        {\url{https://dataforgood.fr/}}
        {}
        {\begin{itemize}
\item Accompagnement d'associations à but non lucratif dans l'analyse de leurs données
\item Développement d'un outil de suivi des bénéficiaires pour une association d'aide alimentaire
\item Animation d'ateliers de sensibilisation aux enjeux éthiques de la donnée
\end{itemize}}

\end{document}